\chapter{State of the art: deployment architectures for a precise standoff}
\label{chapter:sota}

This chapter reports the orientation-phase literature study. It establishes where the existing field leaves a gap that PLASMO can address, provides the basis for the research questions, and establishes which architectures the design trade-off starts from.

\section{Scope and method of the review}
\label{sec:review-method}

The review covers mechanisms whose function is to extend or hold a rigid element to a defined pose: booms, masts, telescoping tubes, deployable-mirror linkages and secondary-mirror supports. Structures whose function is to unfold a large flexible surface, such as mesh and inflatable antennas, membranes and origami arrays, are named and set aside, because holding a rigid point and covering an area are different problems with different measures of success. Approximately 300 papers were catalogued in a breadth scan, and 111 were read and written up recording each mechanism's architecture, its actuation, what defines its deployed position, and the performance reported.

The screening between those two numbers was done with the assistance of a large language model. Each abstract from the breadth scan was passed to the model against a fixed set of questions, namely whether the paper describes a mechanism that extends or holds a rigid element to a defined pose, which family it appears to belong to, and whether it reports deployed performance. Papers the model marked as relevant were then read in full by the author, and the family assignment and the reported performance were taken from that reading rather than from the model. The screening step therefore decides only what is read, never what is concluded.

The purpose of the wider catalogue is coverage rather than measurement. Nothing quantitative is derived from it, and not every catalogued paper is cited. What it supports is the completeness claim of \Cref{sec:families}: that a breadth scan at this size returned no mechanism outside the four kinematic principles and no additional top-level family. The papers that carry the argument are the ones cited in the sections that follow.

The literature contains a large number of performance figures, but few of them can be compared with one another. Collecting the reported packing ratios, accuracies and stiffnesses and placing them on shared axes is possible, but the resulting comparison is not reliable. What a deployed mechanism reports is governed mainly by the quality to which that particular unit was designed, machined, assembled, aligned and measured, and only secondarily by its kinematic architecture. The same architecture built as a laboratory demonstrator and as a flight-qualified unit can differ by significantly in repeatability, and that difference says nothing about the architecture. Measurement conditions widen the gap further: reported figures are variously single-axis or full pose, taken at the end tip or at the mirror, in 1\,g or gravity-compensated, at ambient conditions or in thermal vacuum, and over a handful of deployments or over dozens.

A fair comparison would require each candidate to be rebuilt to the same precision level and tested on the same rig, which is a dedicated experimental programme rather than one chapter of a thesis. The review therefore rests on architecture rather than on reported performance, and draws four kinds of conclusion.

\begin{enumerate}
    \item \textbf{Classification and completeness.} Sorting the architectures into families states what is possible, and the set can be shown to be complete.
    \item \textbf{Properties that follow from the architecture.} Some behaviour comes from the geometry rather than from how well a unit was made: how a structure stows, how many interfaces sit in its load path, whether it is exactly constrained, what holds it in its deployed position, and over what separation range it is viable at all.
    \item \textbf{Absence.} That something has not been addressed is established by searching rather than by measuring, so the objection above does not apply to it. A wide enough literature research can show which combinations of mechanisms have been tried and what hasn't.
    \item \textbf{Feasibility evidence.} A single result can show that an architecture can be built to a given level of performance, which is evidence of feasibility.
\end{enumerate}

Rebuilding every candidate to a common standard is out of reach; modelling them to one is not. The literature review therefore selects which architectures to carry forward, and the comparison between those candidates is generated within the thesis, with each one sized on the same assumptions and the same assumed tolerances.

\section{The mechanism families}
\label{sec:families}

Families are placed on two axes rather than listed ad hoc, so that gaps are visible: the kinematic principle by which the structure changes shape, and the actuation source that drives it. Every mechanism found in the review falls into one of the cells, and empty cells indicate combinations that are not represented in the literature.

% \begin{table}[H]
%     \centering
%     \caption{Classification of deployment architectures by kinematic principle and actuation. A breadth scan of deployable booms, masts and trusses returned no mechanism outside these four principles and no additional top-level family.}
%     \label{tab:families}
%     \small
%     \begin{tabularx}{\textwidth}{p{2.9cm}XXp{2.3cm}p{1.1cm}}
%         \toprule
%         Principle / actuation & Passive (spring or stored strain) & Motor or cable driven & Smart material (SMA, SMP) & Fluid (gas) \\
%         \midrule
%         Rigid-body mechanism (hinges, telescoping, linkage, scissor, truss) & A, B, D, E & B, D, F & Rare (SMA-released hinges) & -- \\
%         Stored strain energy (elastic self-deploying booms) & C (tape-spring, STEM, CTM), coilable-longeron mast & -- & SMP booms (rare) & -- \\
%         Inflation (pressure, optionally rigidised) & -- & -- & -- & G \\
%         Tension and membrane (cables, tensioned surfaces) & H (tensegrity), I (origami membrane) & H (cable-driven) & -- & -- \\
%         \bottomrule
%     \end{tabularx}
% \end{table}

\begin{table}[H]
    \centering
    \caption{Classification of deployment architectures by kinematic principle and actuation. A breadth scan of deployable booms, masts and trusses returned no mechanism outside these four principles and no additional top-level family.}
    \label{tab:families}
    \small
    % \renewcommand{\arraystretch}{1.5}
    \begin{tabularx}{\textwidth}{Xp{2.7cm}p{2.7cm}p{2.7cm}p{0.7cm}}
        \toprule
        Principle / actuation & Passive (spring or stored strain) & Motor or cable driven & Smart material (SMA, SMP) & Fluid \\
        \midrule
        Rigid-body mechanism (hinges, telescoping, linkage, scissor, truss) & A, B, D, E & B, D, F & Rare (SMA-released hinges) & -- \\
        \addlinespace
        Stored strain energy (elastic self-deploying booms) & C (tape-spring, STEM, CTM), coilable-longeron mast & -- & SMP booms (rare) & -- \\
        \addlinespace
        Inflation (pressure, optionally rigidised) & -- & -- & -- & G \\
        Tension and membrane (cables, tensioned surfaces) & H (tensegrity), I (origami membrane) & H (cable-driven) & -- & -- \\
        \bottomrule
    \end{tabularx}
\end{table}

The nine families and their descriptions are listed in \Cref{tab:family-descriptions}.

\begin{table}[H]
    \centering
    \caption{The nine mechanism families, described by the deployed load-path architecture that defines each one.}
    \label{tab:family-descriptions}
    \small
    \begin{tabularx}{\textwidth}{p{0.6cm}p{3.2cm}>{\raggedright\arraybackslash}Xp{1.7cm}}
        \toprule
        & Family & Architecture & Example \\
        \midrule
        \textbf{A} & Articulated booms & Rigid arms on discrete hinges at the bus, folded against it for launch, inward or sideways, and rotated out on deployment. Arms can be connected in series to increase the stand-off distance. The family is bounded by the surface the arms fold against rather than by the direction in which they fold. & \Cref{fig:family-a} \\
        \addlinespace
        \textbf{B} & Telescoping booms & Concentric tubes nested inside one another, each extending from the one before it along the optical axis and locking at the end of its stroke. & \Cref{fig:family-b} \\
        \addlinespace
        \textbf{C} & Coilable and strain-energy booms & A single thin-walled shell (tape-spring, STEM, CTM, coilable longeron), flattened and coiled onto a spool for stowage. On release the flattened section recovers its curvature at a travelling deployment front, and the strain energy released there drives the boom outward and forms the beam behind it, so structure, hinge and actuator are one piece of material. & \Cref{fig:family-c} \\
        \addlinespace
        \textbf{D} & Deployable trusses and masts & Discrete members held by joints or latches, extending from a stowed package into a long triangulated column of repeating bays. Instances deploy as articulated truss units, as screw-driven masts, or as hinged beams unfolding in sequence. & \Cref{fig:family-d} \\
        \addlinespace
        \textbf{E} & Pantograph and scissor & Elements pinned to one another at their crossings and at their ends, forming chains of cross units. One input stroke at the base extends the whole chain, multiplying a small displacement into a large expansion. & \Cref{fig:family-e} \\
        \addlinespace
        \textbf{F} & Driven linkage & A chain of rigid bars with a single degree of freedom, driven at one hinge while the remaining hinges follow passively. Four bars is the common case rather than a defining property; what defines the family is the single driven degree of freedom. The mirror travels along one fixed path and the linkage is locked on arrival, so the drive sets the deployed position. & \Cref{fig:family-f} \\
        \addlinespace
        \textbf{G} & Inflatable and inflatable-rigidi\-sable & A packed envelope inflated by gas to form a pressurised boom, optionally rigidised after deployment. & \Cref{fig:family-g} \\
        \addlinespace
        \textbf{H} & Tensegrity & Struts that do not touch one another, each held in place by a net of pretensioned cables. The structure carries load only once the net is fully tensioned. & \Cref{fig:family-h} \\
        \addlinespace
        \textbf{I} & Origami and tessellated folding & Membranes or panels folded on a repeating crease pattern, unfolding from a compact stack into an extended surface. & \Cref{fig:family-i} \\
        \bottomrule
    \end{tabularx}
\end{table}

% \begin{figure}[H]
%     \centering
%     \includegraphics[width=0.6\textwidth]{figures/family-a.pdf}
%     \caption{Family A: articulated booms. Source: \cite{bouwmeester_enabling_2023}.}
%     \label{fig:family-a}
% \end{figure}

Where architecture and actuation point to different families, the deployed load-path architecture determines the family. For extending columns the distinguishing feature is the presence of joints: discrete members held by joints or latches form a truss column, while a column whose members are themselves the strained material is a strain-energy boom. This criterion follows the literature. Fenci and Currie \cite{fenci_deployable_2017} find that the one distinction recurring across every classification scheme they examine is rigid links against deformable components, and that a unified classification of deployable structures is still missing. A narrower classification arrived at independently supports the same structure: Schwartz et al. \cite{schwartz_high-resolution_2020} group deployable secondary-mirror concepts into telescopic boom, tape spring, truss mast and articulated arm, which correspond to families B, C, D and A above. Closer still is Wang et al. \cite{wang_space_2024}, an independent review that sorts the field on the same two axes used here. Its three morphing categories, pantographic, elastically deformed and inflatable, are the rigid-body, stored-strain and inflation principles of \Cref{tab:families}, and its four driving categories, motor, spring, inflatable and smart material, match the actuation columns one for one. Five schemes have been checked against this morphology in total \cite{fenci_deployable_2017, kiper_deployable_2009, villalba_review_2020, schwartz_high-resolution_2020, wang_space_2024}, sorting variously by kinematics, deployment geometry, failure mechanism and driving principle, and none names a category without a home in the nine families.

% Put this once in the preamble or just above the figure.
\newcommand{\famstub}{\fbox{\rule{0pt}{3.0cm}\rule{4.0cm}{0pt}}}

\begin{figure}[htbp]
    \centering

    \begin{subfigure}[t]{0.31\textwidth}
        \centering
        \famstub % \includegraphics[height=3.0cm]{figures/family-a.pdf}
        \caption{\textbf{A} - Articulated booms \cite{bouwmeester_enabling_2023}}
        \label{fig:family-a}
    \end{subfigure}
    \hfill
    \begin{subfigure}[t]{0.31\textwidth}
        \centering
        \famstub % \includegraphics[height=3.0cm]{figures/family-b.pdf}
        \caption{\textbf{B} Telescoping booms \cite{gooding_novel_2019}}
        \label{fig:family-b}
    \end{subfigure}
    \hfill
    \begin{subfigure}[t]{0.31\textwidth}
        \centering
        \famstub % \includegraphics[height=3.0cm]{figures/family-c.pdf}
        \caption{\textbf{C} Coilable and strain-energy booms \cite{black_deployment_2006}}
        \label{fig:family-c}
    \end{subfigure}

    \vspace{0.8em}

    \begin{subfigure}[t]{0.31\textwidth}
        \centering
        \famstub % \includegraphics[height=3.0cm]{figures/family-d.pdf}
        \caption{\textbf{D} Deployable trusses and masts \cite{kempenaers_deployable_2025}}
        \label{fig:family-d}
    \end{subfigure}
    \hfill
    \begin{subfigure}[t]{0.31\textwidth}
        \centering
        \famstub % \includegraphics[height=3.0cm]{figures/family-e.pdf}
        \caption{\textbf{E} Pantograph and scissor \cite{nagy_design_2021}}
        \label{fig:family-e}
    \end{subfigure}
    \hfill
    \begin{subfigure}[t]{0.31\textwidth}
        \centering
        \famstub % \includegraphics[height=3.0cm]{figures/family-f.pdf}
        \caption{\textbf{F} Driven linkage \cite{ni_four-bar_2019}}
        \label{fig:family-f}
    \end{subfigure}

    \vspace{0.8em}

    \begin{subfigure}[t]{0.31\textwidth}
        \centering
        \famstub % \includegraphics[height=3.0cm]{figures/family-g.pdf}
        \caption{\textbf{G} Inflatable and inflatable-rigidisable \cite{freeland_large_1997}}
        \label{fig:family-g}
    \end{subfigure}
    \hfill
    \begin{subfigure}[t]{0.31\textwidth}
        \centering
        \famstub % \includegraphics[height=3.0cm]{figures/family-h.pdf}
        \caption{\textbf{H} Tensegrity \cite{tibert_deployable_2002}}
        \label{fig:family-h}
    \end{subfigure}
    \hfill
    \begin{subfigure}[t]{0.31\textwidth}
        \centering
        \famstub % \includegraphics[height=3.0cm]{figures/family-i.pdf}
        \caption{\textbf{I} Origami and tessellated folding \cite{XXX_origami}}
        \label{fig:family-i}
    \end{subfigure}

    \caption{One example of each mechanism family. Figures reproduced from the cited sources.}
    \label{fig:families}
\end{figure}

\section{Architectural properties and manufacturing properties}
\label{sec:properties}

Properties set by the architecture, and therefore comparable between concepts:

\begin{itemize}
    \item \textbf{Stowage geometry.} How a structure collapses, and therefore roughly how densely it packs, follows from its kinematics. A boom that coils onto a spool packs differently from an arm that folds against a wall.
    \item \textbf{The number of interfaces in the load path.} Every joint, latch and sliding fit adds a point at which position is handed on. Better manufacturing reduces each contribution but does not remove it, so error accumulation scales with interface count as a property of the architecture. Edeson et al. \cite{edeson_conventional_2010} give the mechanism: partial slip in a joint displaces the structural connection it supports, and the effect is amplified across a structure carrying a number of joints under asymmetric loading.
    \item \textbf{Kinematic determinacy.} Whether the deployed structure is exactly constrained or overdefined is a design property. An overdefined structure with many latches carries internal loads when it heats unevenly, which can cause thermal warping.
    % \item \textbf{The thermal load path.} How many joints the standoff carries, and whether the members can be laid out symmetrically, determines how predictable its thermal expansion is and where a compensating layout could be applied. An athermal design deliberately combines materials of differing expansion so that the net displacement at system level cancels, so mixing materials is the method rather than an obstacle to it; what makes it hard is an unsymmetrical chain of many joints, not the presence of more than one material. Athermalisation is treated here as an optimisation available to most architectures rather than as a discriminator between them, and it is scoped out of the trade-off: earlier work found it either not fully achievable or only partly effective, it enlarges the design because members must be lengthened to give the compensating return path room, and the thermomechanical decoupling of the baffle already keeps the thermal variation seen by the mirror support within a narrow band.
    % \item \textbf{The thermal load path.} How many joints the standoff carries, and whether its members can be laid out symmetrically, sets how predictable its thermal expansion is and where a compensating layout could be applied. An athermal design combines materials of differing expansion so that the net displacement cancels at system level, so using more than one material is the method rather than the obstacle; the difficulty is an unsymmetrical chain of many joints.
    
    \item \textbf{What defines the deployed position.} Three cases behave differently in kind, and this property says most about whether a concept can reach a sub-micron target. Which of the three applies is set at the hinge rather than by the family, so the hinge and locking concept is treated here as a decision in its own right.
    \begin{itemize}
        \item A hard mechanical stop or latch is repeatable, and its accuracy improves with workmanship.
        \item A driven mechanism with feedback places the element and can be corrected.
        \item A structure that finds its own shape from stored elastic energy is subject to hysteresis and, over long stowage, to viscoelastic relaxation. Neither is removed by precision manufacturing, because both follow from deploying by strain energy.
    \end{itemize}
    \item \textbf{The geometric viability envelope.} Folded arms can only be as long as the surface they fold against, whichever direction they fold in, so the architecture becomes unavailable once the mirror separation (and thus the boom length) grows beyond the length or width of the satellite bus, which will be roughly the primary diameter. Extending architectures do not have the same bound.
    \item \textbf{Whether a fine-correction stage can be carried, and where.} Some architectures offer a natural interface at the arm bases or beneath the mirror; others do not.
\end{itemize}

Properties set by manufacture, and therefore not comparable between papers:

\begin{itemize}
    \item \textbf{The repeatability in micrometres.}
    \item \textbf{The drift per kelvin.}
    \item \textbf{The deployed first eigenfrequency in hertz.}
    \item \textbf{The absolute mass.}
    \item \textbf{The achieved packing ratio of any one built article.}
\end{itemize}

These figures are recorded in the underlying notes and are used below only as feasibility evidence.

\section{The TU Delft heritage and comparable missions}
\label{sec:heritage}

% TODO (Jasper): this section is abstract; add illustrations, with correct attribution. Figures may be taken from the source papers.

% PLASMO's deployable optics descend from the TU Delft Deployable Space Telescope programme. Its value for this thesis is that it is a sequence of successive theses on the same architecture, so its open ends, non-conformances and improvement potential are documented rather than inferred.

% Dolkens and Kuiper \cite{dolkens_deployable_2017} set out the end-to-end architecture of a segmented primary with a boom-deployed secondary and active calibration. Van Putten \cite{van_putten_design_2017} designed the driven deployment mechanism for the primary-mirror segments. The secondary-mirror support was then approached twice from opposite directions, which makes it the most informative pair in the line: Krikken \cite{krikken_design_2018} designed it passively, athermal by material and kinematic choice with tensioned ribbons, while Akkerhuis \cite{akkerhuis_deployable_2020} redesigned the same element as a fully active system with piezo actuators and interferometric metrology. Gritter \cite{gritter_assessing_2020} characterised the thermo-mechanical behaviour of the hinge that carries the booms. Every design in the line uses the same hinge arrangement: compliant rolling-contact hinges at the boom ends with a self-locking strip or slotted hinge in the middle, deployed by the strain stored in the hinge itself, with a magnetic latch built and tested alongside as the stiffer alternative. That arrangement follows from a specific argument, which is worth stating because the thesis inherits it: the compliant rolling-contact hinge was developed so that the deployed system would not be overdetermined, on the reasoning that an overdetermined structure warps under uneven heating. A compliant hinge can still move after deployment, but without the backlash and hysteresis characteristic of other hinge types. The argument has never been tested. It has not been simulated whether a locked hinge really would warp more than the accuracy budget allows, nor whether the compliant hinge introduces unwanted behaviour of its own under transient or microvibration loading. This is the open question carried in \Cref{sec:deferred}. Hobijn \cite{hobijn_design_2023} built and tested a secondary-mirror deployment mechanism, which makes it the one hardware repeatability campaign in the line and the source of the practical failure modes reported: hinge hysteresis, collapsing hinges, and trapped air limiting deployment length. Nagy \cite{nagy_design_2021} treated the deployable baffle separately, as a scissor structure. The line arrives at the PLASMO baseline itself: three inward-folding booms carrying the mirror and baffle with a fine adjustment at the root \cite{bouwmeester_enabling_2023}.

PLASMO's deployable optics descend from the TU Delft Deployable Space Telescope programme. Its value for this thesis is that it is a sequence of successive theses on the same architecture, so its open ends, non-conformances and improvement potential are documented rather than inferred. The line runs in two generations: a segmented-primary telescope for visible-light Earth observation, and the smaller thermal-infrared demonstrator that PLASMO inherits directly.

\subsection{Visible-light Deployable Space Telescope}

The first generation begins with Dolkens and Kuiper \cite{dolkens_deployable_2017}, who set out the end-to-end architecture of a segmented primary with a boom-deployed secondary and active calibration, and with Van Putten \cite{van_putten_design_2017}, who designed the driven deployment mechanism for the primary-mirror segments. The support structure that holds the secondary was then approached twice from opposite directions, which makes it the most informative pair in the line: Krikken \cite{krikken_design_2018} designed it passively, athermal by material and kinematic choice with tensioned ribbons, while Akkerhuis \cite{akkerhuis_deployable_2020} redesigned the same element as a fully active system with piezo actuators and interferometric metrology, having first quantified a thermoelastic drift that passive athermalisation could not bring inside the budget. Gritter \cite{gritter_assessing_2020} characterised the thermo-mechanical behaviour of the hinge that carries the booms, and Nagy \cite{nagy_design_2021} treated the deployable baffle as a separate element, a pantographic scissor structure, in work whose spectral band moved to the thermal infrared partway through.

Every design in the line uses the same hinge arrangement: compliant rolling-contact hinges at the boom ends with a self-locking strip or slotted hinge in the middle, deployed by the strain stored in the hinge itself, with a magnetic latch built and tested alongside as the stiffer alternative. That arrangement follows from a specific argument, which is worth stating because the thesis inherits it: the compliant rolling-contact hinge was developed so that the deployed system would not be overdetermined, on the reasoning that an overdetermined structure warps under uneven heating. A compliant hinge can still move after deployment, but without the backlash and hysteresis characteristic of other hinge types. The argument has never been tested. It has not been simulated whether a locked hinge really would warp more than the accuracy budget allows, nor whether the compliant hinge introduces unwanted behaviour of its own under transient or microvibration loading. This is the open question carried in \Cref{sec:deferred}.

\subsection{Thermal-infrared demonstrator}

The second generation changes scale. Hobijn \cite{hobijn_design_2023} designed, built and tested a secondary-mirror deployment mechanism for a \SI{30}{\centi\meter} thermal-infrared demonstrator at \SI{300}{\kilo\meter}, with its requirements scaled down from the visible-light design by the ratio of the wavelengths and the support structure taken as given. It is the one hardware repeatability campaign in the line and the source of the practical failure modes reported: hinge hysteresis, collapsing hinges, and trapped air limiting deployment length. Because the two generations differ in aperture, band and orbit, their reported accuracies describe different problems and are not directly comparable. The line arrives at the PLASMO baseline itself: three inward-folding booms carrying the mirror and baffle with a fine adjustment at the root \cite{bouwmeester_enabling_2023}.

The main difference for PLASMO is the accuracy target. The thermal-infrared line worked to a tolerance set by long wavelengths, whereas PLASMO's visible and short-wave infrared bands tighten it by roughly an order of magnitude. That change is what makes a purely passive solution unlikely, and it moves the passive-versus-active question from a design detail to a central question of the thesis.

\subsection{Other comparable missions}

Outside Delft, the closest comparable works are the CubeSat active-optics line developed at the UK Astronomy Technology Centre \cite{schwartz_laboratory_2018, schwartz_high-resolution_2020, schwartz_active_2021, schwartz_6u_2022, schwartz_high-resolution_2025}, which reaches similar optic sizes on a small platform, folds the secondary as well as the primary segments in its later designs, and is consistently active-optics-led, and JWST \cite{lightsey_james_2012}, which demonstrates the full chain of driven deployment, fine stage and wavefront control, and remains informative despite the difference in scale.


\section{The families assessed}
\label{sec:families-assessed}

Each family is assessed against the properties of \Cref{sec:properties} and against the evidence that it has been built for an optical standoff.The outcome is whether to carry it into the design trade-off or to set it aside.  Where a family set aside is genuinely good at something, that is recorded with it.
% There is no middle shelf: a family that is not modelled cannot be placed beside families that are without reproducing the incomparable comparison \Cref{sec:review-method} rejects.

\subsection*{Carried into the trade-off}

\textbf{A, articulated booms.} Rigid arms on discrete hinges at the bus, folded against it for launch, inward or sideways, and rotating out into a tripod that holds the mirror. The load path contains few interfaces, and the arm bases give a natural interface for a correction stage. How the deployed position is held is an open choice within the family rather than a property of it, and the answer the heritage line reached is described in \Cref{sec:heritage}. This family also has the deepest heritage for this specific task, including the PLASMO baseline \cite{bouwmeester_enabling_2023} and the TU Delft Deployable Space Telescope \cite{krikken_design_2018, akkerhuis_deployable_2020, gritter_assessing_2020, hobijn_design_2023}, so its hinge failure modes are documented: hysteresis, plastic deformation (in slotted-tape hinges and others), loss of stored energy (in compliant hinges), cold welding (in metal-on-metal hinges), backlash, X-Y play and rotation \cite{hobijn_design_2023, liu_thin-walled_2024, villalba_review_2020}. Whether to lock every joint or only the middle one is an open trade in that work, and is the overdefinition question discussed in \Cref{sec:deferred}. The limit of the family is geometric: the arms have to lie against the bus, so it is available only while the mirror separation (and thus the boom length) stays below the length or width of the satellite bus, which will be roughly the primary diameter. Carried as the small-separation candidate.

\textbf{D, deployable trusses and masts.} Discrete-member columns that extend into a long triangulated structure. They are stiffness-efficient over a long reach, which is where family A ceases to be available, and the architecture has been applied to a secondary mirror in at least one recent design \cite{kempenaers_deployable_2025}. The load path contains many joints, so the accuracy case depends on how those joints are latched and preloaded, though a driven screw-type extension has placed a mast tip to within a few micrometres \cite{wang_design_2026}, which shows the architecture is not inherently coarse. Two things count against it at this scale. No published work in the family reports thermal behaviour, so its stability at a sub-micron target is unknown. And truss booms typically need a large deployment canister whose volume can exceed what a small satellite allows \cite{belvin_advanced_2016}. Carried as the large-separation candidate.

\textbf{F, driven linkage.} A chain of rigid bars with one degree of freedom, driven out along a single fixed path and locked on arrival. The instances found are predominantly four-bar, but the bar count is incidental; the single driven degree of freedom is what characterises the family. The deployed position is set by the drive rather than by a hard stop or by stored strain, but with one degree of freedom the mechanism places the mirror along that path only, and locking removes that: Ni et al. \cite{ni_four-bar_2019} report zero degrees of freedom once locked, at about \SI{1}{\milli\meter} and 0.05 degrees, so a separate correction stage is still needed behind it. The same architecture carries the JWST secondary mirror at a much larger scale, a four-bar of four composite tubes with one driven and four passive hinges, and there the requirement on the linkage is only to land the mirror inside the capture range of the alignment loop \cite{lightsey_james_2012}. Stowage does not shape the family, since the swept arc has to be kept clear, and no instance in it reports a packing ratio. Carried as the driven-deployment option.

\textbf{B, telescoping booms.} Concentric tubes that slide out along the optical axis and lock. The family reaches beyond the bus dimension and is light. However, the sliding fit is the position-defining feature, and sliding fits are subject to stiction, jamming and thermal binding. The published instances are coarse by design \cite{reveles_development_2015}, and one pneumatic demonstrator stalled part-way through deployment \cite{badea_proof_2025}. Carried only as a structural carrier for a driven extension or a fine stage, not as a precision architecture in its own right.


\subsection*{Set aside}

\textbf{C, coilable and strain-energy booms.} A single thin-walled shell coiled for stowage, in which the stored elastic energy provides the actuation and the recovered cross-section forms the beam. Structure, hinge and actuator are one piece of material, so the family stows densely and can weigh very little, and it has flight heritage \cite{belvin_advanced_2016, wang_space_2024, liu_thin-walled_2024}. The objection is architectural, the deployed shape is the material's own equilibrium, so hysteresis is intrinsic \cite{black_deployment_2006}, and prolonged stowage, which lasts months to years, evolves the viscoelastic and plastic properties of the resin, reducing the stored strain energy and degrading deployment and shape recovery \cite{liu_thin-walled_2024, thaker_tape_2024}. Set aside. It could serve as a packing reference, but a family that is not modelled cannot be set beside families that are without reproducing exactly the unfair comparison \Cref{sec:review-method} rejects, and its intrinsic hysteresis rules it out for this standoff in any case. What it is good at is recorded here: it stows more densely and weighs less than any other rigid-member architecture, the inflatable family excepted, and it reaches long distances for very little mass.

\textbf{E, pantograph and scissor.} Chains of cross-linked elements that convert a small stroke into a large expansion. Expansion ratio is high, but the load path passes through a large number of pin joints, so tolerance accumulation is inherent to the concept, and the same interface count brings low stiffness and a high part count. Set aside for the mirror standoff. What it is good at is recorded here: it is proven and appropriate for the deployable baffle, where inaccuracy is tolerated \cite{nagy_design_2021}, and it can be stiff in particular directions once extended.

\textbf{G, inflatable and inflatable-rigidisable.} Can pack very small, at minimal stored strain energy, and reaches long standoffs \cite{schenk_review_2014}, but the unsolved problem is PLASMO's binding constraint, holding a precise and stable pose. The most complete review of inflatable boom packing and rigidisation concludes that the accuracy required by reflectors and optical components exceeds what inflatables can readily achieve, and that the coefficient of thermal expansion of rigidisable materials is often too great for such missions \cite{schenk_review_2014}. The instances found in this family are surface expanders, inflatable reflectors and antenna arrays whose measure of success is how much area they unfold and how closely the surface matches its intended shape \cite{freeland_large_1997, huang_development_2001, kamaliya_inflatable_2023}: within the scope defined in \Cref{sec:review-method}, no reviewed instance performs an optical standoff, and none reports a deployed position for a rigid element.

\textbf{H, tensegrity.} Discontinuous struts in a pretensioned cable net. Stiffness exists only once the structure is fully tensioned and is very low in bending, which is a property of the concept \cite{tibert_deployable_2002}.

\textbf{I, origami and tessellated folding.} A surface technology for sunshields, baffles and apertures. It unfolds an area rather than holding a rigid point at a standoff, so it is out of scope for the M2 structure and relevant only if the separate baffle work is revisited.

% \section{Agreement and silence in the literature}
% \label{sec:agreement}

% Three positions recur across the literature independently of how well any individual article was built.

% \begin{itemize}
%     \item Sub-micron performance is reached with active correction. Every reviewed system that achieves it does so behind a correction loop, and none claims to reach it passively; Villalba et al. \cite{villalba_review_2020} report this as the consensus of the deployable-optics literature. What has to be held here is the placement of M2 with respect to M1, a tolerance that follows from the instrument, the wavelength and the resolution to be achieved, and that falls at the micron to sub-micron level. The question for this thesis is therefore how far passive design reaches against that placement budget, and what correcting the remainder costs, not whether correction is needed at all.
%     \item Baffle deployment is usually treated as a separate design problem from the secondary-mirror mechanism, as it is in the TU Delft Deployable Space Telescope and in the PLASMO baseline itself \cite{nagy_design_2021, bouwmeester_enabling_2023}, but not universally, as there are three projects that design the deployable baffle as part of the payload, two of them carrying it on the mirror structure, and they differ in how far they take it.
%     \begin{itemize}
%         \item The SSTL deployable telescope makes the deployed structure the enclosure, three concentric barrels with the secondary mirror carried on the upper one, and states that its baffling deploys with the telescope; that system is asserted rather than shown, with no geometry, stray-light budget or thermal analysis \cite{gooding_novel_2019, shore_new_2019}.
%         \item The CubeSat Cassegrain of Kempenaers and Vandepitte \cite{kempenaers_deployable_2025} is the one worked case: three baffles deploy with the structure, the secondary baffle mounted on the M2 assembly, sized against stray light by non-sequential ray tracing, but the thermomechanical effect of carrying them is not analysed.
%         \item The CubeSat active-optics line at the UK Astronomy Technology Centre deploys a baffle alongside the mirrors while keeping it a separate subsystem at concept level, and in its later designs moves stray-light control into the optical train \cite{schwartz_high-resolution_2020, schwartz_6u_2022, schwartz_high-resolution_2025}. 
%     \end{itemize} 
%     Integration is therefore demonstrated in concept, dimensioned optically once, and nowhere assessed for what it does to the stability of the mirror it carries. The reason usually given does not settle it either, since differing accuracy requirements would only mean designing to the tighter one. The binding reason is thermal: the baffle intercepts direct sunlight and albedo, so a mirror support sharing that structure inherits the resulting thermomechanical dynamics, which is what makes the tighter requirement hard to hold.
%     \item The correction stage is designed as its own object. Fine-positioning mechanisms for secondary mirrors form a distinct and mature line of work, separate from the deployment mechanisms that carry them.
% \end{itemize}

% The following positions do not occur in the researched literature.

% \begin{itemize}
% % TODO (Jasper, 27 Aug 2026) RESEARCH NEEDED before this bullet stands.
% % "If you are specific enough you will eventually find no literature. The question is how relevant the combination is. Is the VLEO environment really very different from LEO? If so, in what respect?" He assumes examples in the 20-30 cm class do exist. VLEO itself is new, so it may well be the interesting angle, but the text has to name WHAT is new: the thermal environment, or the atomic oxygen and drag.
% % Action: research the LEO vs VLEO difference and rewrite this bullet around the mechanism behind the gap.

%     \item A packing ratio at the level of the deployment mechanism, stated on a common basis, is almost never given. Reports give either the whole telescope or a single boom, and rarely say which. Whole-telescope figures are reasonable in themselves, since the system level is what a mission cares about, but they mix the mechanism with everything else on the spacecraft and so cannot be used to compare architectures.
    
%     \item No deployment mechanism is designed around an off-axis aperture, in which the support cannot be symmetric about the beam. Off-axis instruments in the literature hold their secondary in a fixed, non-deploying mount. The nearest deployable case is off-axis only in its field: the TU Delft telescope uses an annular-field Korsch design whose pupil stays centred and obscured, so its secondary hangs on a symmetric three-boom tripod \cite{dolkens_deployable_2017}.

%     \item The optical penalty of the deployed structure itself is not quantified. Whether members sit in the optical path, how many of them there are, and how the structure behaves for stray light are all consequences of the architecture, and they admit a first-order comparison: three struts obstruct less than four, and a spider in the field of view is unavoidable for some families and not for others.

%     % \item The volume the support structure itself demands between mirror and baffle is not reported. It matters because the baffle has to clear the structure: the more room the mechanism takes there, the larger the baffle becomes. While launch volume is a design driver in this satellite size range, at this altitude that propagates further than launch volume. Once the baffle is deployed it is the largest ram-facing surface on the satellite, so its cross-section sets the drag force while its length sets the offset between centre of pressure and centre of mass, and therefore the aerodynamic torque the attitude control has to absorb continuously \cite{aslanov_vleo_2023}. With the centre of mass near the bus, the force grows roughly with baffle length and the torque with its square. A mechanism that forces a wider or longer enclosure is paid for in drag, in control authority and in orbit lifetime, not only in fairing space.
    
%     % \item The thermal case of a deployed support inside its own enclosure is not analysed anywhere. Where the enclosure is modelled in detail, the boom set is inherited and fixed \cite{hu_thermo-optical_2022}; where the booms are modelled, there is no enclosure at all, and deployable insulating sleeves are proposed as future work \cite{de_zanet_predicted_2023}. Since a baffle surrounding the optical assembly makes its inner wall the mechanism's radiative environment, this coupled case is the one that governs, and neither paper poses it. Separately, every controlled thermal test in the reviewed work is a uniform soak, so the gradient case that produces tilt and decentre is never applied.

%     % \item The atomic-oxygen exposure of an enclosed deployed structure is not quantified. Published work tests flat, unstrained coupons at ram conditions, and the flight failure mode is breach of the protective layer at edges, creases and clamped regions \cite{goto_changes_2023, smith_complete_2021}. Inside a baffle the mechanism sees mostly re-emitted, thermally accommodated atoms, at energies well below those at which any published erosion yield was measured, while the entrance region retains direct exposure at full energy. How far the enclosure has to extend beyond the mirror for that to hold is set in the literature by stray-light requirements alone, and no source treats it as an exposure parameter.

%     \item The volume the support structure itself demands between mirror and baffle is not reported. It matters because the baffle has to clear the structure: the more room the mechanism takes there, the larger the baffle becomes. While launch volume is a design driver in this satellite size range, at this altitude that propagates further than launch volume. At \SI{300}{\kilo\meter} the atmospheric density is one to two orders of magnitude above the 500 to \SI{600}{\kilo\meter} band in which the deployable-optics literature sits, and the drag force rises by more than two hundredfold between \SI{600}{\kilo\meter} and \SI{300}{\kilo\meter} \cite{crisp_benefits_2020}. Once the baffle is deployed it is the largest ram-facing surface on the satellite, so its cross-section sets the drag force while its length sets the offset between centre of pressure and centre of mass, and therefore the aerodynamic torque the attitude control has to absorb continuously \cite{aslanov_vleo_2023}. With the centre of mass near the bus, the force grows roughly with baffle length and the torque with its square. A mechanism that forces a wider or longer enclosure is paid for in drag, in control authority and in orbit lifetime, not only in fairing space. The same enclosure at \SI{600}{\kilo\meter} carries almost none of that cost, which is why no reviewed source reports the figure.

%     \item The thermal case of a deployed support inside its own enclosure is not analysed anywhere. Where the enclosure is modelled in detail, the boom set is inherited and fixed \cite{hu_thermo-optical_2022}; where the booms are modelled, there is no enclosure at all, and deployable insulating sleeves are proposed as future work \cite{de_zanet_predicted_2023}. Since a baffle surrounding the optical assembly makes its inner wall the mechanism's radiative environment, this coupled case is the one that governs, and neither paper poses it. The altitude changes which terms drive that inner wall rather than how severe the cycle is: at \SI{300}{\kilo\meter} against \SI{600}{\kilo\meter} the Earth infrared and albedo fluxes on a nadir-facing surface rise by about 9\%, the orbital period shortens from 96.5 to \SI{90.4}{\minute}, and the direct solar term is unchanged, so the planetary contribution is a larger share of a faster drive while the eclipse case warms by roughly \SI{4.5}{\kelvin} and the orbital temperature swing narrows by about \SI{2}{\kelvin}. That is an argument for posing the environment as a time-dependent profile at a stated beta angle rather than as a constant worst case \cite{gonzalez-barcena_selection_2022}, not for expecting a harsher cycle. Separately, every controlled thermal test in the reviewed work is a uniform soak, so the gradient case that produces tilt and decentre is never applied.

%     \item The atomic-oxygen exposure of an enclosed deployed structure is not quantified. Published work tests flat, unstrained coupons at ram conditions, and the flight failure mode is breach of the protective layer at edges, creases and clamped regions \cite{goto_changes_2023, smith_complete_2021}. At \SI{300}{\kilo\meter} the atomic-oxygen density is at least an order of magnitude above the 500 to \SI{700}{\kilo\meter} band in which the reviewed exposures and the deployable-optics work sit, and molecular nitrogen reaches 10\% or more of the residual atmosphere, arriving at about \SI{8.8}{\electronvolt} against \SI{5.1}{\electronvolt} for the oxygen \cite{goto_changes_2023}. Inside a baffle the mechanism sees mostly re-emitted, thermally accommodated atoms, at energies well below those at which any published erosion yield was measured, while the entrance region retains direct exposure at full energy. How far the enclosure has to extend beyond the mirror for that to hold is set in the literature by stray-light requirements alone, and no source treats it as an exposure parameter, so the shielding factor is assumed rather than established against a flux an order of magnitude higher than any published exposure.

% \end{itemize}

% \section{The gap and the novelty claim}
% \label{sec:gap}

% % No published work designs or analyses a deployable secondary-mirror structure for a 20 to \SI{30}{\centi\meter} aperture in very low Earth orbit, at a sub-micron accuracy and stability target.

% % The nearest work covers part of it. The deployable-telescope heritage works either to a coarser accuracy target or at a much larger aperture. The CubeSat active-optics line \cite{schwartz_laboratory_2018, schwartz_high-resolution_2020, schwartz_active_2021, schwartz_6u_2022, schwartz_high-resolution_2025} comes closest, and its later designs fold a secondary mirror as well as the primary segments at a comparable aperture, but it gets its accuracy from active phasing of the optics rather than from the structure, and it works to a \SI{500}{\kilo\meter} baseline without looking at the thermal environment lower down. The long-reach mast work is structural rather than optical and reports no thermal behaviour.

% No published work establishes how far passive design reaches against the placement and stability budget of a deployable secondary mirror held inside its own stray-light enclosure, nor what a fine stage is then left to absorb.

% The aperture class and the \SI{300}{\kilo\meter} orbit are the parameters of that question rather than the gap itself. The aperture fixes the standoff, and with it the lever arm on every angular error. The altitude fixes the size of the enclosure through drag, and the material environment through atomic-oxygen flux, and it changes which terms drive the thermal case without making the cycle harsher.

% The nearest work covers part of it. The deployable-telescope heritage works either to a coarser accuracy target or at a much larger aperture. The CubeSat active-optics line \cite{schwartz_laboratory_2018, schwartz_high-resolution_2020, schwartz_active_2021, schwartz_6u_2022, schwartz_high-resolution_2025} comes closest on aperture, and its later designs fold a secondary mirror as well as the primary segments, but it obtains its accuracy from active phasing of the optics rather than from the structure, and it works to a \SI{500}{\kilo\meter} baseline, so neither the enclosure coupling nor the material environment at \SI{300}{\kilo\meter} enters its analysis. The long-reach mast work is structural rather than optical and reports no thermal behaviour.

% Six things follow from this.

% \begin{itemize}
%     \item \textbf{A like for like comparison.} \Cref{sec:review-method} explains why the published figures cannot be set against each other. So the thesis builds its own comparison instead: every carried architecture modelled on one set of assumptions, at one assumed precision and material class, so that the differences come from the architecture rather than from who built the hardware and how well.
%     \item \textbf{A packing ratio at the level that matters.} Published reports give whole-telescope volumes, or the length ratio of a single boom. Neither isolates what a designer choosing between deployment architectures needs to know. This thesis works the figure at the level of the M2 mechanism: the volumetric ratio of the mechanism stowed to the same mechanism deployed, bounded where it mounts to the rest of the spacecraft and running to its furthest extent, so that it excludes both the individual boom below it and the satellite and its other deployables above it. It is worked out the same way for every architecture, reported as a range between the minimum envelope and the practical enclosing box, and accompanied by a separate figure for the volume the baffle is forced to occupy.
    
%     % \item \textbf{Thermal stability at \SI{300}{\kilo\meter}.} In-orbit stability under orbital thermal cycling is reported for very few mechanisms, and for none in this regime. The thesis runs the thermal cases, feeds them into a thermoelastic analysis, and carries the result into the error budget.
    
%     \item \textbf{Thermal stability of the support inside its enclosure.} In-orbit stability under orbital thermal cycling is reported for very few mechanisms, and the coupled case for none: the one detailed enclosure model inherits a fixed boom set \cite{hu_thermo-optical_2022}, and the one orbital thermal model of a deployable secondary support has no enclosure at all and is not validated against any measurement of a deployable \cite{de_zanet_predicted_2023}. Since the baffle's inner wall is what the mechanism radiates to, that coupled case is the one this thesis runs. The altitude changes which terms drive it rather than how severe it is, so the contribution lies in how the case is posed: a time-dependent flux profile at a stated beta angle, chosen against the structure's own thermal response time \cite{gonzalez-barcena_selection_2022}, applied to more than one architecture on common assumptions, and including the gradient case rather than a uniform soak. The result feeds the thermoelastic analysis and the error budget, and is reported as a band rather than a point.
    
%     \item \textbf{The switch point.} At what M1 to M2 separation do folding architectures stop being the better answer and extending mechanisms (like telescoping and coilable booms) take over, and is that boundary set by the ratio to the mirror diameter or by an absolute distance? The literature does not address it. Since the optical design is still open, an answer that holds across the range is worth more than a design fixed to one separation.
    
%     \item \textbf{How far passive design gets.} The literature agrees that active correction is needed. What it does not say is how much error is left once the passive measures are exhausted, and the closest predecessor assumes the answer rather than testing it \cite{hobijn_design_2023}. The thesis states that residual, and the stroke and resolution a fine stage would need to absorb it.
    
%     % TODO (Jasper, 27 Aug 2026) DECISION NEEDED, revisit in mid-September.
%     % "Certainly interesting given the current status. But this alternative can generate a lot of work, so keep something in reserve. If in a few weeks it turns out to be on-axis anyway, that is an opportunity to descope." His estimate is that the mission goes on-axis: the central obstruction costs only ~10-11%, an off-axis M2 needs a much larger baffle, and both mirrors would have to go freeform, so strong radiometric or optical arguments would be needed to defend it. Off-axis is also awkward to carry into the parametric models because it comes in many kinds. Action: keep this bullet for now; drop it, and sub-question 2.2 with it, once the instrument design settles on-axis.
%     \item \textbf{An off-axis aperture.} No deployment mechanism in the literature is built for one, and off-axis instruments hold their secondary on fixed mounts. The thesis asks whether a single concept can serve both layouts, which matters because that choice is not yet fixed, even though the mission currently leans on-axis.
% \end{itemize}

\section{Agreement and silence in the literature}
\label{sec:agreement}

Three positions recur across the literature independently of how well any individual article was built.

\begin{itemize}
    \item Sub-micron performance is reached with active correction. Every reviewed system that achieves it does so behind a correction loop, and none claims to reach it passively; Villalba et al. \cite{villalba_review_2020} report this as the consensus of the deployable-optics literature. What has to be held here is the placement of M2 with respect to M1, a tolerance that follows from the instrument, the wavelength and the resolution to be achieved, and that falls at the micron to sub-micron level. That budget is not the wavefront-error budget which governs a segmented primary, and the two should not be read as one.
    \item Baffle deployment is usually treated as a separate design problem from the secondary-mirror mechanism, as it is in the TU Delft Deployable Space Telescope and in the PLASMO baseline itself \cite{nagy_design_2021, bouwmeester_enabling_2023}, but not universally, as there are three projects that design the deployable baffle as part of the payload, two of them carrying it on the mirror structure, and they differ in how far they take it.
    \begin{itemize}
        \item The SSTL deployable telescope makes the deployed structure the enclosure, three concentric barrels with the secondary mirror carried on the upper one, and states that its baffling deploys with the telescope, with no geometry, stray-light budget or thermal analysis given \cite{gooding_novel_2019, shore_new_2019}.
        \item The CubeSat Cassegrain of Kempenaers and Vandepitte \cite{kempenaers_deployable_2025} is the one worked case: three baffles deploy with the structure, the secondary baffle mounted on the M2 assembly, sized against stray light by non-sequential ray tracing, but the thermomechanical effect of carrying them is not analysed.
        \item The CubeSat active-optics line at the UK Astronomy Technology Centre deploys a baffle alongside the mirrors while keeping it a separate subsystem at concept level, and in its later designs moves stray-light control into the optical train \cite{schwartz_high-resolution_2020, schwartz_6u_2022, schwartz_high-resolution_2025}.
    \end{itemize}
    Integration is therefore demonstrated in concept and dimensioned optically once. The reason usually given for avoiding it does not settle it on its own, since differing accuracy requirements would only mean designing to the tighter one. The binding reason is thermal: the baffle intercepts direct sunlight and albedo, so a mirror support sharing that structure inherits the resulting thermomechanical dynamics, which is what makes the tighter requirement hard to hold.
    \item The correction stage is designed as its own object. Fine-positioning mechanisms for secondary mirrors form a distinct and mature line of work, separate from the deployment mechanisms that carry them.
\end{itemize}

The following positions do not occur in the researched literature.

\begin{itemize}
    \item A packing ratio at the level of the deployment mechanism, stated on a common basis, is almost never given. Reports give either the whole telescope or a single boom, and rarely say which. Whole-telescope figures are reasonable in themselves, since the system level is what a mission cares about, but they mix the mechanism with everything else on the spacecraft and so cannot be used to compare architectures.

    \item No deployment mechanism is designed around an off-axis aperture, in which the support cannot be symmetric about the beam. Off-axis instruments in the literature hold their secondary in a fixed, non-deploying mount. The nearest deployable case is off-axis only in its field: the TU Delft telescope uses an annular-field Korsch design whose pupil stays centred and obscured, so its secondary hangs on a symmetric three-boom tripod \cite{dolkens_deployable_2017}.

    \item The optical penalty of the deployed structure itself is not quantified. Whether members sit in the optical path, how many of them there are, and how the structure behaves for stray light are all consequences of the architecture, and they admit a first-order comparison: three struts obstruct less than four, and a spider in the field of view is unavoidable for some families and not for others.

    \item The volume the support structure itself demands between mirror and baffle is not reported. It matters because the baffle has to clear the structure: the more room the mechanism takes there, the larger the baffle becomes. While launch volume is a design driver in this satellite size range, at this altitude that propagates further than launch volume. At \SI{300}{\kilo\meter} the atmospheric density is one to two orders of magnitude above the 500 to \SI{600}{\kilo\meter} band in which the deployable-optics literature sits, and the drag force rises by more than two hundredfold between \SI{600}{\kilo\meter} and \SI{300}{\kilo\meter} \cite{crisp_benefits_2020}. Once the baffle is deployed it is the largest ram-facing surface on the satellite, so its cross-section sets the drag force while its length sets the offset between centre of pressure and centre of mass, and therefore the aerodynamic torque the attitude control has to absorb continuously \cite{aslanov_vleo_2023}. With the centre of mass near the bus, the force grows roughly with baffle length and the torque with its square. A mechanism that forces a wider or longer enclosure is paid for in drag, in control authority and in orbit lifetime, not only in fairing space. The same enclosure at \SI{600}{\kilo\meter} carries almost none of that cost, and the reviewed work sits at that altitude and above.

    \item The thermal case of a deployed support inside its own enclosure is not analysed anywhere. Where the enclosure is modelled in detail, the boom set is inherited and fixed \cite{hu_thermo-optical_2022}; where the booms are modelled, there is no enclosure at all, and deployable insulating sleeves are proposed as future work \cite{de_zanet_predicted_2023}. Nor do the three concepts that carry the enclosure on the deployed structure assess what doing so costs the stability of the mirror they carry. The enclosure is not neutral: it shields the structure from direct solar and albedo flux while itself heating and radiating to it, which is both the reason integration is usually avoided and the reason the coupled case cannot be left out of a separated design either. The altitude changes which terms drive that inner wall rather than how severe the cycle is: at \SI{300}{\kilo\meter} against \SI{600}{\kilo\meter} the Earth infrared and albedo fluxes on a nadir-facing surface rise by about 9\%, the orbital period shortens from 96.5 to \SI{90.4}{\minute}, and the direct solar term is unchanged, so the planetary contribution is a larger share of a faster drive while the eclipse case warms by roughly \SI{4.5}{\kelvin} and the orbital temperature swing narrows by about \SI{2}{\kelvin}. That is an argument for posing the environment as a time-dependent profile at a stated beta angle rather than as a constant worst case \cite{gonzalez-barcena_selection_2022}, not for expecting a harsher cycle. Separately, every controlled thermal test in the reviewed work is a uniform soak, so the gradient case that produces tilt and decentre is never applied.

    \item The atomic-oxygen exposure of an enclosed deployed structure is not quantified. Published work tests flat, unstrained coupons at ram conditions, and the flight failure mode is breach of the protective layer at edges, creases and clamped regions \cite{goto_changes_2023, smith_complete_2021}. At \SI{300}{\kilo\meter} the atomic-oxygen density is at least an order of magnitude above that of the 500 to \SI{600}{\kilo\meter} band in which the deployable-optics work sits, the factor being measured in flight by Goto et al. against the 500 to \SI{700}{\kilo\meter} band \cite{goto_changes_2023}, and molecular nitrogen reaches 10\% or more of the residual atmosphere, arriving at about \SI{8.8}{\electronvolt} against \SI{5.1}{\electronvolt} for the oxygen. Inside a baffle the mechanism sees mostly re-emitted, thermally accommodated atoms, at energies well below those at which any published erosion yield was measured, while the entrance region retains direct exposure at full energy. How far the enclosure has to extend beyond the mirror for that to hold is set in the literature by stray-light requirements alone, and no source treats it as an exposure parameter, so the shielding factor is assumed rather than established against a flux an order of magnitude higher than any published exposure.
\end{itemize}

\section{The gap and the novelty claim}
\label{sec:gap}

No published work establishes how far passive design reaches against the placement and stability budget of a deployable secondary mirror held inside its own stray-light enclosure, nor what a fine stage is then left to absorb.

The aperture class and the \SI{300}{\kilo\meter} orbit are the parameters of that question rather than the gap itself. The aperture fixes the standoff, and with it the lever arm on every angular error. The altitude fixes the size of the enclosure through drag, changes which terms drive the thermal case without making the cycle harsher, and sets the material environment through atomic-oxygen flux.

The nearest work covers part of it. The deployable-telescope heritage works either to a coarser accuracy target or at a much larger aperture. The CubeSat active-optics line \cite{schwartz_laboratory_2018, schwartz_high-resolution_2020, schwartz_active_2021, schwartz_6u_2022, schwartz_high-resolution_2025} comes closest on aperture, and its later designs fold a secondary mirror as well as the primary segments, but it obtains its accuracy from active phasing of the optics rather than from the structure, and it works to a \SI{500}{\kilo\meter} baseline, so neither the enclosure coupling nor the material environment at \SI{300}{\kilo\meter} enters its analysis. The long-reach mast work is structural rather than optical and reports no thermal behaviour.

Six things follow from this.

\begin{itemize}
    \item \textbf{A like for like comparison.} \Cref{sec:review-method} explains why the published figures cannot be set against each other. So the thesis builds its own comparison instead: every carried architecture modelled on one set of assumptions, at one assumed precision and material class, so that the differences come from the architecture rather than from who built the hardware and how well.

    \item \textbf{A packing ratio at the level that matters.} Published reports give whole-telescope volumes, or the length ratio of a single boom. Neither isolates what a designer choosing between deployment architectures needs to know. This thesis works the figure at the level of the M2 mechanism: the volumetric ratio of the mechanism stowed to the same mechanism deployed, bounded where it mounts to the rest of the spacecraft and running to its furthest extent, so that it excludes both the individual boom below it and the satellite and its other deployables above it. It is worked out the same way for every architecture, reported as a range between the minimum envelope and the practical enclosing box, and accompanied by a separate figure for the volume the baffle is forced to occupy.

    \item \textbf{Thermal stability of the support inside its enclosure.} In-orbit stability under orbital thermal cycling is reported for very few mechanisms, and the coupled case for none: the one detailed enclosure model inherits a fixed boom set \cite{hu_thermo-optical_2022}, and the one orbital thermal model of a deployable secondary support has no enclosure at all and is not validated against any measurement of a deployable \cite{de_zanet_predicted_2023}. Since the baffle's inner wall is what the mechanism radiates to, that coupled case is the one this thesis runs. The altitude changes which terms drive it rather than how severe it is, so the contribution lies in how the case is posed: a time-dependent flux profile at a stated beta angle, chosen against the structure's own thermal response time \cite{gonzalez-barcena_selection_2022}, applied to more than one architecture on common assumptions, and including the gradient case rather than a uniform soak. The result feeds the thermoelastic analysis and the error budget, and is reported as a band rather than a point.

    \item \textbf{The switch point.} At what M1 to M2 separation do folding architectures stop being the better answer and extending mechanisms (like telescoping and coilable booms) take over, and is that boundary set by the ratio to the mirror diameter or by an absolute distance? The literature does not address it. Since the optical design is still open, an answer that holds across the range is worth more than a design fixed to one separation.

    \item \textbf{How far passive design gets.} The closest predecessor assumes the residual rather than establishing it \cite{hobijn_design_2023}, so the figure a correction stage would have to be sized against does not exist anywhere. The thesis produces it: the error remaining once the passive measures are exhausted, and the stroke and resolution a fine stage needs to absorb it.

    % TODO (Jasper, 27 Aug 2026) DECISION NEEDED, revisit in mid-September.
    % "Certainly interesting given the current status. But this alternative can generate a lot of work, so keep something in reserve. If in a few weeks it turns out to be on-axis anyway, that is an opportunity to descope." His estimate is that the mission goes on-axis: the central obstruction costs only ~10-11%, an off-axis M2 needs a much larger baffle, and both mirrors would have to go freeform, so strong radiometric or optical arguments would be needed to defend it. Off-axis is also awkward to carry into the parametric models because it comes in many kinds. Action: keep this bullet for now; drop it, and sub-question 2.2 with it, once the instrument design settles on-axis. If dropped, "Six things follow from this" becomes five.
    \item \textbf{An off-axis aperture.} No deployment mechanism in the literature is built for one, and off-axis instruments hold their secondary on fixed mounts. The thesis asks whether a single concept can serve both layouts, which matters because that choice is not yet fixed, even though the mission currently leans on-axis.
\end{itemize}

\section{Research questions}
\label{sec:questions}

% TODO (Jasper, 27 Aug 2026) DECISION NEEDED on the framing of the whole question set. "You could make it a microsatellite telescope for Earth observation, with PLASMO as the design case. That way the parameter ranges of interest are not necessarily set by the current state of the design, and the question becomes rather more scientific, while keeping a strong case-driven design focus." Same point from the research side in the meeting: the thesis should serve the project AND anyone interested in deployable optics, which argues for taking parameters as variable inputs from the start rather than running a sensitivity analysis afterwards. If adopted, this rewords the main question and sub-question 2.

% TODO (open since 26 Aug 2026) RATIFICATION NEEDED, not raised at the review.
% Sub-question 1 as locked on 16 July named a shortlist: "(articulated fold-in with a 3-DOF fine stage; screw-driven truss/mast; telescoping carrier; gated inflatable)". That parenthetical was removed on 26 August because naming candidates in a research question pre-commits the trade-off. That is a change to a question Jasper agreed, and it did not come up on 27 August. Put it to him on 10 September.

% \subsection*{Main research question}

% \begin{quote}
% \textbf{How can a deployable structure for the PLASMO instrument's secondary mirror be designed to meet the required deployment accuracy, structural stiffness and in-orbit stability while maximising the packing ratio, given that the optical configuration (aperture within \SI{20}{\centi\meter} to \SI{30}{\centi\meter}, on-axis versus off-axis, and the mirror sizes and distances) is not yet fixed?}
% \end{quote}

% \subsection*{Sub-questions}

% \begin{enumerate}
%     \item \textbf{Concepts and mechanisms.} Which deployment-mechanism families and hinge or locking concepts are feasible for the secondary-mirror structure, and how do the architectures carried forward compare on volumetric packing ratio, bare deployment accuracy, in-orbit stability, stiffness, mass, complexity and required active correction?
%     \item \textbf{Validity across the open optical space.} How can the deployment design remain valid across the open optical solution space, in particular:
%     \begin{enumerate}
%         \item At what M1 to M2 separation does the preferred architecture switch from folding to extending, is that switch governed by a relative measure (a ratio to the mirror diameter) or by an absolute separation value?
%         \item Can a single concept span both regimes, on-axis and off-axis?
%     \end{enumerate}
%     \item \textbf{The passive-active split.} To what extent is active control needed to compensate thermo-mechanical effects on deployment accuracy and in-orbit stability, and\slash or which passive measures (material choice, kinematics, or one-time calibration) can reduce or remove that need? Deployment accuracy is a one-time event, whereas stability is continuous, so the two are not necessarily answered by the same measure.
% \end{enumerate}

% \subsection*{Main research question}

% \begin{quote}
% \textbf{How can a deployable structure for the PLASMO instrument's secondary mirror, operating at very low altitude and inside its own deployable stray-light enclosure, be designed to meet the required deployment accuracy, structural stiffness and in-orbit stability while maximising the packing ratio, given that the optical configuration (aperture within \SI{20}{\centi\meter} to \SI{30}{\centi\meter}, on-axis versus off-axis, and the mirror sizes and distances) is not yet fixed?}
% \end{quote}

% \subsection*{Sub-questions}

% \begin{enumerate}
%     \item \textbf{Concepts and mechanisms.} Which deployment-mechanism families and hinge or locking concepts are feasible for the secondary-mirror structure, and how do the architectures carried forward compare on volumetric packing ratio, the enclosure volume each one forces, bare deployment accuracy, in-orbit stability, stiffness, mass, complexity and required active correction?
%     \item \textbf{Validity across the open optical space.} How can the deployment design remain valid across the open optical solution space, in particular:
%     \begin{enumerate}
%         \item At what M1 to M2 separation does the preferred architecture switch from folding to extending, is that switch governed by a relative measure (a ratio to the mirror diameter) or by an absolute separation value?
%         \item Can a single concept span both regimes, on-axis and off-axis?
%     \end{enumerate}
%     \item \textbf{The passive-active split.} To what extent is active control needed to compensate the thermo-mechanical effects of the orbital environment at \SI{300}{\kilo\meter}, taken as a time-dependent profile inside the deployed enclosure, on deployment accuracy and in-orbit stability, and\slash or which passive measures (material choice under the exposure at that altitude, kinematics, or one-time calibration) can reduce or remove that need? Deployment accuracy is a one-time event, whereas stability is continuous, so the two are not necessarily answered by the same measure.
% \end{enumerate}


\subsection*{Main research question}

\begin{mainquestion}
    \item \label{rq:main} \textbf{How can a deployable structure for the PLASMO instrument's secondary mirror, operating at very low altitude and inside its own deployable stray-light enclosure, be designed to meet the required deployment accuracy, structural stiffness and in-orbit stability while maximising the packing ratio, given that the optical configuration (aperture within \SI{20}{\centi\meter} to \SI{30}{\centi\meter}, on-axis versus off-axis, and the mirror sizes and distances) is not yet fixed?}
\end{mainquestion}

\subsection*{Sub-questions}

\begin{researchquestions}
    \item \label{rq:concepts} \textbf{Concepts and mechanisms.} Which deployment-mechanism families and hinge or locking concepts are feasible for the secondary-mirror structure, and how do the architectures carried forward compare on volumetric packing ratio, the enclosure volume each one forces, bare deployment accuracy, in-orbit stability, stiffness, mass, complexity and required active correction?
    \item \label{rq:validity} \textbf{Validity across the open optical space.} How can the deployment design remain valid across the open optical solution space, in particular:
    \begin{researchquestions}
        \item \label{rq:switch} At what M1 to M2 separation does the preferred architecture switch from folding to extending, is that switch governed by a relative measure (a ratio to the mirror diameter) or by an absolute separation value?
        \item \label{rq:axis} Can a single concept span both regimes, on-axis and off-axis?
    \end{researchquestions}
    \item \label{rq:passive} \textbf{The passive-active split.} To what extent is active control needed to compensate the thermo-mechanical effects of the orbital environment at \SI{300}{\kilo\meter}, taken as a time-dependent profile inside the deployed enclosure, on deployment accuracy and in-orbit stability, and\slash or which passive measures (material choice under the exposure at that altitude, kinematics, or one-time calibration) can reduce or remove that need? Deployment accuracy is a one-time event, whereas stability is continuous, so the two are not necessarily answered by the same measure.
\end{researchquestions}